% ---
% id: receiver/rem:rankone
% kind: rem
% title: "Ports of rank greater than one"
% summary: >-
%   No result of this contract assumes one optical mode per detector port: a
%   port collecting several modes can have a Gram of rank above one. The
%   one-mode-per-port description holds exactly under the rank criterion of
%   [[thm:representation-independence]] (iv).
% prominence: standard
% status: accepted
% uses: [thm:representation-independence]
% mentions: [thm:click-povm, thm:coarse-graining, thm:multiclick-bound, thm:emitted-bound, thm:exact-and-evaluated, lem:deficiency, receiver/detector-imperfections/thm:efficiency, receiver/detector-imperfections/thm:dark-counts, receiver/component-catalogue/ex:bb84, receiver/detector-imperfections/rem:effrank, req:one-mode-per-port]
% assumes: []
% refs: []
% provenance: newly-derived
% lean:
% ---

\begin{remark}\label{rem:rankone}
Example~\ref{ex:bb84} ends with four rank-one Grams, which is why the textbook
$4\times2$ BB84 instrument is legitimate \emph{there}. Nothing in this
contract assumes one mode per port anywhere else: the click-pattern POVM,
the theorems built on it and the decomposition of the \tref{loss deficiency},
among them Theorems~\ref{thm:click-povm}, \ref{thm:coarse-graining},
\ref{thm:multiclick-bound}, \ref{thm:emitted-bound} and
\ref{thm:exact-and-evaluated} and Lemma~\ref{lem:deficiency} of the core
and Theorems~\ref{thm:efficiency} and~\ref{thm:dark-counts} of its
submodules, are proved for arbitrary $\{A_k\}$, and a port that collects
several modes, such as a multimode fiber onto one active area or two spatial
modes onto one detector, can have $\rk A_k>1$, and every result above holds
when it does. The criterion for a
one-mode-per-port description is Theorem~\ref{thm:representation-independence}\,(iv); it is decidable
from the emitted Grams, and Remark~\ref{rem:effrank} exhibits the one
operation in this contract, apart from declaring a zero efficiency, that can
change its answer.
Requirement~\ref{req:one-mode-per-port} states the rule the engine follows.
\end{remark}
